\section{\uppercase{Discussion}}
\label{sec:discussion}

\textbf{Interpretation.} The overall observed order, with pose and coarse structure settling first while fine detail settles last, matches what is known about unconditional diffusion models~\cite{choi2022p2,rissanen2023ihdm}. Since in virtual try-on, the model uses a garment image that already contains all of the detail it needs to generate, this was not self-evident. Nonetheless, the models we evaluated still follow the natural frequency order imposed by the noise and the garment condition only has an effect on \emph{how early} the detail can be inpainted. This is where the models appear to vary. A possible explanation lies in how they inject the garment. OOTDiffusion, for example, processes the clean garment with a separate outfitting UNet and fuses those features into every self-attention layer of the denoising UNet. This way, detailed garment features are available from the first step, which fits its early rise in texture and pattern. On the other hand, CatVTON concatenates the garment to the person in the same input, causing the garment's detail to have to be propagated by the same network that is still resolving the noise. Finally, IDM-VTON also uses a separate garment UNet, but its timing is closer to CatVTON than to OOTDiffusion, therefore it appears that the garment encoder alone does not completely explain the difference. Additional possible factors that our study cannot separate include the base model (SDXL versus SD~1.5), the IP-Adapter, the different guidance scales and the different native resolutions. It is worth noting that although OOTDiffusion commits early, it also produces the most stable trajectories, so its early commitment does not come at the cost of more corrections later in the process.
\\ \\
\textbf{Implications.} Given that for all three models pose and structure settle within the first few steps, we conclude that the early steps fix the layout and the later steps primarily add the garment's appearance. This makes the second half of the trajectory, where texture, pattern and colour settle in \catvton and \idm, the place to intervene when print fidelity is the problem, for example with a stronger garment guidance in that part of the schedule. This result also supports a handover strategy like the one used in LPH-VTON~\cite{liu2026lphvton}. Given our measurements, the handover should happen around step 15 to 20 of 50 for CatVTON and IDM-VTON. Finally, the long, relatively flat tail of changes following the large modifications suggests that the number of sampling steps could be reduced, although the colour axis shows that the colours appear to be continuously refined until the end.
\\ \\
\textbf{Limitations.} The estimate $\hat{x}_0$ is merely an estimate of the final image, not the image itself. Accordingly, at high noise levels it is close to an average over many possible images. Therefore, a part of the coarse-to-fine behaviour is caused by the estimate rather than by the generation. This affects all models in the same way, however, thus the differences between models remain meaningful. Except for stability, our curves are measured against the final frame of the same run. Consequently, they show when a property reaches its final state, but not whether it was already decided at that point or merely becomes visible. Experiments that re-sample the noise from an intermediate step on would be needed to separate these two options. Furthermore, the axes are proxies. Colour is not independent of structure and pattern is not independent of texture. All results are on upper-body garments from a single dataset, with 48 pairs, one seed and one sampler per model. Two of the models were run with a sampler other than their default one. The standard scores in Table~\ref{tab:benchmark} are taken from the literature and were not measured on our runs. Finally, although all frames are scored at the same resolution, the models generate at their own native resolution ($384 \times 512$ for CatVTON, $768 \times 1024$ for the other two), which may affect how quickly fine detail appears.
\\ \\
\textbf{Future work.} To increase the confidence of these observations, natural next steps include evaluating more models and datasets, such as Dress Code~\cite{morelli2022dresscode} for lower-body garments and dresses. Additionally, restarting the generation from an intermediate step several times with fresh noise would show whether a property is already decided at that step, in which case it no longer varies across the restarts, or only becomes visible later.
